% !TEX root = ../main.tex
% =====================================================================
% Appendices
%   A  Sample Calculations
%   B  Sample Simulation
%   C  Prototype
%   D  Preliminary Experiments
%   E  Candidate Solutions
% Appendix chapters are lettered; their equations are unnumbered so they
% do not enter the List of Equations.
% =====================================================================
\clearpage
\setcounter{section}{0}
\renewcommand{\thesection}{\Alph{section}}
\renewcommand{\theHsection}{appendix.\Alph{section}}
\titleformat{\section}[display]{\normalfont\bfseries\centering}{Appendix \Alph{section}}{0pt}{\MakeUppercase}
\addcontentsline{toc}{section}{APPENDICES}

% =====================================================================
\section[APPENDIX A.\ SAMPLE CALCULATIONS]{Sample Calculations}
\label{app:calc}

This appendix works through the formulas of Chapter~\ref{sec:method} with numbers small enough to check by hand. Sections~\ref{app:calc-grid} to \ref{app:calc-fail} follow one toy utterance through the stopping policy; Section~\ref{app:calc-ltt} computes the calibration p-values and the sample-size limit. In the toy example each ``output'' has only three numbers, the threshold is $\lambda=0.50$ and the SNR gate is $\theta=20$~dB. The values illustrate the computation and are not experimental results.

\subsection{Teacher Grid and Cost of Each Candidate}
\label{app:calc-grid}

With \Cref{eq:grid}, $t_4 = 1-\cos(\tfrac{\pi}{2}\cdot\tfrac{4}{16}) = 1-\cos(0.3927) = 1-0.9239 = 0.0761$, and the other grid points follow in the same way. With \Cref{eq:cost}, stopping at $k=4$ costs $2(4+1)+3 = 13$ forward passes. Table~\ref{tab:app-grid} lists all grid points up to the reference.

\begin{table}[!htbp]
\centering\small
\caption{Teacher grid points (\Cref{eq:grid}) and the cost of each candidate stopping step (\Cref{eq:cost}).}
\label{tab:app-grid}
\begin{tabular}{@{}cccccc@{}}
\toprule
Step $k$ & $t_k$ & Role & Teacher passes & Cost $C$ & Saving vs.\ 21 \\
\midrule
0--3 & 0 to 0.0431 & Euler steps only & --- & --- & --- \\
4 & 0.0761 & candidate ($s=0.07$) & 10 & 13 & 8 (38.1\%) \\
5 & 0.1181 & candidate ($s=0.10$) & 12 & 15 & 6 (28.6\%) \\
6 & 0.1685 & candidate ($s=0.15$) & 14 & 17 & 4 (19.0\%) \\
7 & 0.2270 & not a candidate & --- & --- & --- \\
8 & 0.2929 & reference ($s=0.25$) & 18 & 21 & 0 \\
\bottomrule
\end{tabular}
\end{table}

\subsection{Prompt SNR Estimate}

Suppose the prompt has ten frames with energies 0.010, 0.012, 0.00001, 0.0002, 0.011, 0.009, 0.0003, 0.013, 0.010 and 0.0002. The quietest 10\% is one frame, so $\bar E_{\text{bottom}} = 0.00001$. The loudest half is 0.010, 0.010, 0.011, 0.012 and 0.013, so $\bar E_{\text{top}} = 0.0112$. By \Cref{eq:snr},
\[
\hat S = 10\log_{10}\frac{0.0112}{0.00001} = 10\log_{10}1120 = 30.5\ \text{dB}.
\]
Since $30.5\ge\theta=20$, the SNR gate is open and early stopping is allowed (Algorithm~\ref{alg:handover}, line~1).

\subsection{Endpoint Estimate}

At step 5 ($t_5=0.1181$), take one entry of the state, $x_{t_5}=0.8$, with guided velocity $v=1.5$ (\Cref{eq:guided}). By \Cref{eq:endpoint} the endpoint estimate of that entry is
\[
\hat x_1(t_5) = 0.8 + (1-0.1181)\times1.5 = 0.8 + 0.8819\times1.5 = 2.12.
\]
The velocity is computed anyway for the Euler step, so this estimate costs no forward pass (Proposition~\ref{prop:zero}).

\subsection{Endpoint-Stabilization Score and Stopping Decision}

Let the three-number endpoint estimates at steps 3, 4 and 5 be $[0.5,\,0.2,\,0.9]$, $[2.0,\,1.0,\,2.0]$ and $[2.2,\,1.0,\,1.6]$. With \Cref{eq:stab}:
\begin{itemize}
\item Step 4 (first candidate): the change is $[1.5,\,0.8,\,1.1]$ with norm $\sqrt{2.25+0.64+1.21}=\sqrt{4.10}=2.025$; the estimate's norm is $\sqrt{4+1+4}=3.000$; so $d_4 = 2.025/3.000 = 0.675 > \lambda = 0.50$. By \Cref{eq:policy} the teacher continues (Algorithm~\ref{alg:handover}, line~7).
\item Step 5: the change is $[0.2,\,0,\,-0.4]$ with norm $\sqrt{0.20}=0.447$; the estimate's norm is $\sqrt{4.84+1+2.56}=\sqrt{8.40}=2.898$; so $d_5 = 0.447/2.898 = 0.154 \le 0.50$, and $\hat S\ge\theta$. The teacher stops at step 5 (lines~4--5).
\end{itemize}
By \Cref{eq:cost} the utterance costs $2(5+1)+3 = 15$ forward passes instead of 21, a saving of $6/21 = 28.6\%$. Had $\hat S$ been 17.8~dB, the gate would have kept the teacher running to the reference step however settled the estimates looked, and the cost would have been 21.

\subsection{Added-Failure Check}
\label{app:calc-fail}

With \Cref{eq:failure} and $c=12.5\%$, $m_W=1$ point, $m_S=0.02$:
\begin{itemize}
\item \emph{No failure.} Reference output: WER 0\%, SIM 0.80. Candidate output: WER 0\%, SIM 0.79. The WER limit is $\max(12.5\%,\,0\%+1\%)=12.5\%$ and the SIM limit is $0.80-0.02=0.78$; both are respected, so $f=0$.
\item \emph{Failure.} Reference output: WER 15.38\% (2 errors in 13 words). Candidate output: WER 23.08\% (3 errors). The WER limit is $\max(12.5\%,\,15.38\%+1\%)=16.38\%$, which is exceeded, so $f=1$ regardless of SIM.
\end{itemize}

\subsection{Calibration p-Values and Sample Size}
\label{app:calc-ltt}

Take $n=100$ calibration inputs, $\alpha=0.05$ and $\delta=0.10$, and suppose the thresholds of the grid produce the failure counts in Table~\ref{tab:app-ltt}. For brevity the example uses a four-value grid; the grid of the study is given in Section~\ref{sec:calibration}. Each p-value follows from \Cref{eq:pvalue}; for example, with one failure,
\[
\begin{aligned}
p &= \P\big(\mathrm{Bin}(100,0.05)\le1\big) = 0.95^{100} + 100\times0.05\times0.95^{99}\\
  &= 0.0059 + 0.0312 = 0.0371.
\end{aligned}
\]
Algorithm~\ref{alg:ltt} certifies $\lambda=0.40$ and $\lambda=0.50$, stops at $\lambda=0.60$, and returns $\hat\lambda=0.50$. By \Cref{eq:guarantee}, with probability at least 90\%, this threshold adds failures to at most 5\% of new inputs from the same population.

\begin{table}[!htbp]
\centering\small
\caption{Fixed-sequence calibration with $n=100$, $\alpha=0.05$, $\delta=0.10$ (illustrative failure counts).}
\label{tab:app-ltt}
\begin{tabular}{@{}cccccl@{}}
\toprule
Order $g$ & $\lambda_g$ & Failures $n\hat R_g$ & $p_g$ (\Cref{eq:pvalue}) & $p_g\le0.10$? & Action \\
\midrule
1 & 0.40 & 0 & 0.0059 & yes & certify, continue \\
2 & 0.50 & 1 & 0.0371 & yes & certify, continue \\
3 & 0.60 & 3 & 0.2578 & no & stop \\
4 & 0.70 & 6 & not tested & --- & --- \\
\bottomrule
\end{tabular}
\end{table}

\emph{Smallest calibration set with zero failures.} With no failures, \Cref{eq:pvalue} gives $p=(1-\alpha)^n$. Certifying $\alpha=0.05$ at $\delta=0.10$ needs $0.95^n\le0.10$, i.e.\ $n\ge\ln0.10/\ln0.95 = 44.9$. Indeed $0.95^{44}=0.105$ fails and $0.95^{45}=0.099$ passes, so at least 45 independent inputs are needed. With 36 inputs, $p=0.95^{36}=0.158>0.10$, which is why the earlier cross-validation (Appendix~\ref{app:prelim}) could not certify this tolerance.

% =====================================================================
\section[APPENDIX B.\ SAMPLE SIMULATION]{Sample Simulation}
\label{app:sim}

This appendix checks, on simulated data, whether the calibration of Section~\ref{sec:obj3} keeps its promise $\P(R\le\alpha)\ge1-\delta$, and compares it with a naive rule that picks the most aggressive threshold whose observed failure rate is at most $\alpha$. No speech is synthesized; the data are random numbers with the structure of Chapter~\ref{sec:method}.

\subsection{Setup}

Each simulated utterance has a prompt SNR, stabilization scores $d_4,d_5,d_6$ at the three candidate steps, and an added-failure indicator for each candidate. A quarter of the utterances are noisy, with $\hat S\sim\mathcal N(16,2^2)$~dB; the rest are clean, with $\hat S\sim\mathcal N(33,4^2)$~dB. An utterance difficulty $h\sim\mathrm{LogNormal}(0,0.45^2)$ drives the scores, which shrink at later steps:
\[
d_k = h\, e^{-0.35(k-4)}\, u_k,\qquad u_k\sim\mathrm{LogNormal}(-0.55,\,0.25^2),\qquad k\in\{4,5,6\}.
\]
A candidate fails with a probability that grows with its score and is higher on noisy prompts:
\[
\P\big(f(x,k)=1\big) = \sigma\big(-4.4 + 3.2\,d_k + 1.6\,\mathbb 1[\text{noisy}] - 0.45\,(k-4)\big),
\]
where $\sigma$ is the logistic function; the reference never fails, as in Definition~\ref{def:failure}. The policy of \Cref{eq:policy} uses $\theta=20$~dB and the grid $\lambda\in\{0.20,0.25,\dots,1.20\}$, with costs from \Cref{eq:cost}. The true risk of every threshold is computed on 400{,}000 simulated utterances (Figure~\ref{fig:app-sim}). Calibration is then repeated 1{,}000 times for each calibration size $n$, with $\alpha=0.05$ and $\delta=0.10$, and the true risk of each returned threshold is recorded. The script, \texttt{sim/ltt\_sim.py}, accompanies this manuscript.

\begin{figure}[!htbp]
\centering
\begin{tikzpicture}
\begin{axis}[
  width=0.82\textwidth, height=5.6cm,
  xlabel={Threshold $\lambda$}, ylabel={True risk $R(\pi_\lambda)$},
  xmin=0.18, xmax=1.22, ymin=0, ymax=0.08,
  yticklabel style={/pgf/number format/fixed, /pgf/number format/precision=2},
  ticklabel style={font=\small}, label style={font=\small},
  axis y line*=left, axis x line*=bottom, grid=major, grid style={gray!20},
]
\addplot[thick, mark=*, mark size=1.2pt, color=newLine] coordinates {
(0.20,0.0021)(0.25,0.0042)(0.30,0.0067)(0.35,0.0096)(0.40,0.0130)(0.45,0.0163)(0.50,0.0202)(0.55,0.0243)(0.60,0.0281)(0.65,0.0318)(0.70,0.0357)(0.75,0.0399)(0.80,0.0439)(0.85,0.0478)(0.90,0.0512)(0.95,0.0549)(1.00,0.0585)(1.05,0.0617)(1.10,0.0651)(1.15,0.0682)(1.20,0.0713)};
\addplot[dashed, thick, color=offLine] coordinates {(0.18,0.05)(1.22,0.05)};
\node[font=\scriptsize, color=offLine, anchor=south west] at (axis cs:0.20,0.05) {tolerance $\alpha=0.05$};
\addplot[only marks, mark=diamond*, mark size=3pt, color=black] coordinates {(0.85,0.0478)};
\node[font=\scriptsize, anchor=north west] at (axis cs:0.86,0.046) {largest safe $\lambda=0.85$};
\end{axis}
\end{tikzpicture}
\caption{True added-failure risk of the stopping policy for each threshold in the simulated population. Larger thresholds stop earlier and fail more often; thresholds up to 0.85 are safe at $\alpha=0.05$.}
\label{fig:app-sim}
\end{figure}

\subsection{Results}

Table~\ref{tab:app-sim} summarizes the 1{,}000 repetitions. Learn-Then-Test met its promise at every calibration size: the returned threshold had true risk at most 5\% in 92.0\% to 98.8\% of repetitions, above the required 90\%. The naive rule did not: because it picks the threshold that happens to look best on one sample, its threshold exceeded the tolerance in more than half of the repetitions, even though its average risk was close to 5\%.

The price of the guarantee is computation. With 50 calibration inputs, Learn-Then-Test was conservative, certified no threshold in about one repetition in ten, and used 18.2 forward passes on average; with 400 inputs it used 15.9, close to the 15.4 of the largest safe threshold. This is the trade-off that the calibration-size analysis of Objective~4 (Section~\ref{sec:obj4}) will measure on real speech.

\begin{table}[!htbp]
\centering\small
\caption{Simulated calibration over 1{,}000 repetitions per size ($\alpha=0.05$, $\delta=0.10$). ``Reference only'' means no threshold was certified. The largest safe threshold, 0.85, uses 15.4 passes on average; the reference uses 21.}
\label{tab:app-sim}
\begin{adjustbox}{max width=\linewidth}
\begin{tabular}{@{}cl ccccc@{}}
\toprule
$n$ & Method & $\P(R\le\alpha)$ & Mean $R$ & Mean passes & Reference only & Median $\hat\lambda$ \\
\midrule
50  & Learn-Then-Test & 0.967 & 0.013 & 18.20 & 10.1\% & 0.35 \\
    & Naive           & 0.427 & 0.052 & 15.54 & 0.0\%  & 0.95 \\
100 & Learn-Then-Test & 0.988 & 0.015 & 17.70 & 2.2\%  & 0.40 \\
    & Naive           & 0.399 & 0.054 & 15.44 & 0.0\%  & 0.95 \\
200 & Learn-Then-Test & 0.967 & 0.027 & 16.46 & 0.0\%  & 0.55 \\
    & Naive           & 0.420 & 0.052 & 15.44 & 0.0\%  & 0.90 \\
400 & Learn-Then-Test & 0.920 & 0.035 & 15.92 & 0.0\%  & 0.70 \\
    & Naive           & 0.469 & 0.051 & 15.43 & 0.0\%  & 0.90 \\
\bottomrule
\end{tabular}
\end{adjustbox}
\end{table}

The simulation assumes independent inputs from one distribution, exactly the condition of Proposition~\ref{prop:ltt}. It therefore shows that the procedure is implemented and behaves as the theory says, not that the condition holds for real speech; that is what the held-out and repeated-split checks of Section~\ref{sec:obj3} will test.

% =====================================================================
\section[APPENDIX C.\ PROTOTYPE]{Prototype}
\label{app:prototype}

The study's prototype is an inference method rather than a user-facing application: the released DMOSpeech~2 sampler with the stopping controller of Section~\ref{sec:architecture} added. Figure~\ref{fig:glance} shows what the prototype changes for a user. The same networks produce the speech; only the point where the teacher hands over to the student moves, separately for each input, and that point never moves later than in the baseline.

The prototype will be packaged as one module shared by all experiment scripts (Section~\ref{sec:setup}). It takes the prompt audio, the prompt text, the target text, the calibrated threshold $\hat\lambda$ and the SNR gate $\theta$, and returns the waveform together with the chosen stopping step, the number of forward passes used, the prompt SNR estimate and the trace of stabilization scores. Before any experiment the module will pass the three checks of Section~\ref{sec:setup}: identical outputs for repeated runs with the same seed, exact reproduction of the reference when the controller falls back, and forward-pass counts equal to \Cref{eq:cost}.

The three boxes at the bottom of Figure~\ref{fig:glance} correspond to what the study must establish about the handover point: how it is decided (Objectives~1 and~2), how safe the decision is (Objective~3, Proposition~\ref{prop:ltt}), and that it never costs more than the baseline (Proposition~\ref{prop:fallback}).

\begin{landscape}
\begin{figure}[p]
\centering
\resizebox{\linewidth}{!}{%
\begin{tikzpicture}[x=1cm,y=1cm,
  frozen/.style={draw=frozenLine, fill=frozenFill, rounded corners=3pt, minimum height=9mm, align=center, font=\small, text=black!75},
  newbox/.style={draw=newLine, fill=white, line width=0.9pt, rounded corners=3pt, align=left, font=\small, inner sep=5pt},
  offbox/.style={draw=offLine, fill=white, line width=0.8pt, rounded corners=3pt, align=left, font=\small, inner sep=5pt},
  lbl/.style={font=\scriptsize, text=black!70, align=center}]
\def\xo{3.4}
\def\u{0.5}
\tikzset{
  tbar/.style={draw=frozenLine, fill=frozenLine!45, line width=0.6pt},
  sbar/.style={draw=frozenLine, fill=frozenFill, line width=0.6pt},
  savebar/.style={draw=newLine, pattern=north east lines, pattern color=newLine!70, line width=0.6pt},
  rowlbl/.style={anchor=east, font=\small, align=right, text=black!80},
  inbar/.style={font=\scriptsize, text=black!80},
}
% (a) baseline
\node[anchor=west, font=\small\bfseries] at (-0.3,1.25) {(a) Baseline: one fixed handover for every input};
\foreach \y/\name in {0/{Input A\\[-2pt]\scriptsize clean, easy}, -0.85/{Input B\\[-2pt]\scriptsize moderate}, -1.7/{Input C\\[-2pt]\scriptsize noisy, hard}}{
  \node[rowlbl] at (\xo-0.15,\y) {\name};
  \draw[tbar] (\xo,\y-0.28) rectangle ({\xo+18*\u},\y+0.28);
  \node[inbar] at ({\xo+9*\u},\y) {teacher \ 18 passes};
  \draw[sbar] ({\xo+18*\u},\y-0.28) rectangle ({\xo+21*\u},\y+0.28);
  \node[inbar] at ({\xo+19.5*\u},\y) {student 3};
}
\draw[frozenLine, dashed, line width=0.9pt] ({\xo+18*\u},0.55) -- ({\xo+18*\u},-2.05);
\node[lbl, above] at ({\xo+18*\u},0.55) {fixed handover $t_h=0.293$};
\node[frozen, anchor=west, text width=88mm, align=left, font=\small] at (15.2,-0.85)
  {Every input pays for the same teacher work (21 passes in total), whether it needs it or not.\\[2pt]
   Easy inputs spend teacher steps that change little in the final speech.};
% (b) proposed
\def\yb{-3.4}
\node[anchor=west, font=\small\bfseries, text=newLine] at (-0.3,\yb) {(b) Proposed: the handover is decided per input, with a guarantee};
\begin{scope}[on background layer]
  \fill[newFill] ({\xo+10*\u},\yb-0.75) rectangle ({\xo+18*\u},\yb-3.35);
  \draw[newLine, line width=1.2pt, rounded corners=2pt] ({\xo+10*\u},\yb-0.75) rectangle ({\xo+18*\u},\yb-3.35);
\end{scope}
\node[font=\scriptsize\bfseries, text=newLine, align=center, anchor=south] at ({\xo+14*\u},\yb-0.72)
  {CONTRIBUTION: where the handover happens};
\foreach \yy/\name/\tp/\saved/\note in
   {1.25/{Input A\\[-2pt]\scriptsize clean, easy}/10/8/{saved 8 passes (38\%)},
    2.10/{Input B\\[-2pt]\scriptsize moderate}/14/4/{saved 4 passes (19\%)},
    2.95/{Input C\\[-2pt]\scriptsize noisy, hard}/18/0/{0 saved (= baseline)}}{
  \pgfmathsetmacro{\yc}{\yb-\yy}
  \node[rowlbl] at (\xo-0.15,\yc) {\name};
  \draw[tbar] (\xo,\yc-0.28) rectangle ({\xo+\tp*\u},\yc+0.28);
  \node[inbar] at ({\xo+\tp*\u/2},\yc) {teacher \tp};
  \draw[sbar] ({\xo+\tp*\u},\yc-0.28) rectangle ({\xo+(\tp+3)*\u},\yc+0.28);
  \node[inbar] at ({\xo+(\tp+1.5)*\u},\yc) {stud.\ 3};
  \ifnum\saved>0
    \draw[savebar] ({\xo+(\tp+3)*\u},\yc-0.28) rectangle ({\xo+21*\u},\yc+0.28);
  \fi
  \node[draw=white, fill=newLine, diamond, inner sep=1.8pt, line width=0.6pt] at ({\xo+\tp*\u},\yc) {};
  \node[anchor=west, font=\scriptsize\bfseries, text=newLine] at ({\xo+21*\u+0.15},\yc) {\note};
}
\foreach \tp/\t in {10/0.076, 12/0.118, 14/0.169, 18/0.293}{
  \draw[newLine, line width=0.8pt] ({\xo+\tp*\u},\yb-3.35) -- ({\xo+\tp*\u},\yb-3.5);
  \node[font=\scriptsize, text=newLine] at ({\xo+\tp*\u},\yb-3.7) {$\t$};
}
\node[font=\scriptsize, text=newLine, anchor=east] at ({\xo+10*\u-0.5},\yb-3.7) {candidate handover times $t_h$:};
\begin{scope}[shift={(\xo,\yb-4.45)}]
  \draw[tbar] (0,0) rectangle (0.6,0.3); \node[anchor=west, font=\scriptsize] at (0.7,0.15) {teacher (frozen, 2 passes/step)};
  \draw[sbar] (5.8,0) rectangle (6.4,0.3); \node[anchor=west, font=\scriptsize] at (6.5,0.15) {student (frozen, 1 pass/step)};
  \draw[savebar] (11.6,0) rectangle (12.2,0.3); \node[anchor=west, font=\scriptsize] at (12.3,0.15) {computation saved};
  \node[fill=newLine, diamond, inner sep=1.8pt] at (16.2,0.15) {}; \node[anchor=west, font=\scriptsize] at (16.4,0.15) {handover chosen by the stopping rule};
\end{scope}
\node[newbox, anchor=north west, text width=78mm, minimum height=17mm] at (-0.3,\yb-4.9)
  {\textbf{\textcolor{newLine}{1. Decide}} \ (\textit{when} to hand over)\\[2pt]
   Stop once the teacher's guess of the final speech stops changing, $d_i\le\hat\lambda$.
   Uses values already computed: \textbf{0 extra passes}.};
\node[offbox, anchor=north west, text width=78mm, minimum height=17mm] at (8.9,\yb-4.9)
  {\textbf{\textcolor{offLine}{2. Guarantee}} \ (\textit{how safe} it is)\\[2pt]
   $\hat\lambda$ is calibrated offline so that added quality failures stay $\le\alpha$
   with probability $\ge 1-\delta$.};
\node[newbox, anchor=north west, text width=78mm, minimum height=17mm] at (18.1,\yb-4.9)
  {\textbf{\textcolor{newLine}{3. Fall back}} \ (\textit{never worse} in cost)\\[2pt]
   Never hands over later than the baseline; if no threshold is certified, the method \emph{is} the baseline.};
\end{tikzpicture}}
\caption{The prototype at a glance. Bar length is computation in forward passes. (a)~The baseline hands over at the same point for every input, so every input costs 21 passes. (b)~The prototype moves the handover inside the highlighted region, separately for each input: settled, easy inputs hand over early and save computation (hatched); hard inputs keep the full teacher run and cost exactly what the baseline costs. Costs are exact (Proposition~\ref{prop:cost}); which input stops where is illustrative.}
\label{fig:glance}
\end{figure}
\end{landscape}


% =====================================================================
\section[APPENDIX D.\ PRELIMINARY EXPERIMENTS]{Preliminary Experiments}
\label{app:prelim}

The team's earlier experiments are reported here as preliminary evidence. They motivated the architecture of Section~\ref{sec:architecture} and fixed several of its choices, but they are not evidence for the calibrated method. The fixed-switching analysis and the cross-validated stabilization policy were run on the Tesla T4 GPU with the evaluation models of Section~\ref{sec:measures}. The student-restart experiment used 12 speakers from Seed-TTS test-en \parencite{anastassiou2024seedtts} under four conditions: a clean prompt, the target duration shortened and lengthened by 15\%, and a prompt with white Gaussian noise added at 15~dB, giving 48 instances. It ran on a central processing unit (CPU) with one seed per instance, measured WER with Whisper-base \parencite{radford2023whisper} and SIM with a lightweight speaker encoder against the prompt, and is therefore not directly comparable with the measures of Section~\ref{sec:measures}. Table~\ref{tab:prelim} summarizes what each experiment found and what it changed in the design.

\subsection{Endpoint Cushioning}

On 96 LibriSpeech pairs, the team tested accepting only part of the student's first correction, $E = T + w(S-T)$, with a tanh-based weight $w = 1 - 0.26\tanh(\|S-T\|_2/\|T\|_2)$, a fixed weight, a cosine-derived weight and a displacement cap. Recomputed from the benchmark log, tanh did not differ significantly from direct transfer in peak spectral flux (mean $-0.41$, paired $t$-test $p=0.31$; 50 better, 39 worse, 7 tied), nor from a matched fixed weight ($p=0.19$). WER changed on only 9 of 96 pairs. The tanh weight ranged only from 0.845 to 0.891, so it behaved almost as a constant, and the cap never activated. The 96 pairs came from 7 speakers, one contributing 41. The student's first correction influences the output, but partial acceptance gave no detectable benefit, so this direction was not pursued as the contribution.

\subsection{Fixed Switching Times}

This analysis used all 500 speakers of the three splits (Section~\ref{sec:data}), each under the clean and the 15~dB noisy prompt condition (1{,}000 instances), synthesized on the Tesla T4 GPU with the student restart fixed, evaluated with the models of Section~\ref{sec:measures}, and scored with the added-failure rule of Definition~\ref{def:failure} against the reference switching time $s=0.25$ (Step~8; \nolinkurl{code/analysis/problem_validation_t4.py}). With one switching time for every input, the failure rate rose as the switch moved earlier, and almost all failures came from noisy prompts (Table~\ref{tab:fixed-fail}); the same ordering held on each split separately. The cheapest switching time that did not fail varied between inputs (Table~\ref{tab:fixed-oracle}), and the 179 instances that needed the reference were all noisy. A per-input choice made in hindsight would therefore use 14.85 forward passes on average instead of 21 ($-29\%$), whereas no single earlier switching time is safe for every input.

\begin{table}[!htbp]
\centering\small
\caption{Added failures of each fixed switching time on 500 speakers under two prompt conditions (Tesla T4), when every input uses the same switching time. Saving is the share of the reference's 21 forward passes that the switching time avoids; the last three columns count added failures (Definition~\ref{def:failure}). An earlier fixed switch saves more but fails more often, almost entirely on noisy prompts.}
\label{tab:fixed-fail}
\setlength{\tabcolsep}{5pt}
\begin{tabular}{@{}l c r r r r r@{}}
\toprule
\textbf{Switching time} & \textbf{Step} & \textbf{Passes} & \textbf{Saving} & \textbf{Clean (500)} & \textbf{Noisy (500)} & \textbf{All (1{,}000)} \\
\midrule
$s=0.07$ & 4 & 13 & 38.1\% & 12 (2.4\%) & 322 (64.4\%) & 334 (33.4\%) \\
$s=0.10$ & 5 & 15 & 28.6\% & 9 (1.8\%) & 269 (53.8\%) & 278 (27.8\%) \\
$s=0.15$ & 6 & 17 & 19.0\% & 5 (1.0\%) & 247 (49.4\%) & 252 (25.2\%) \\
$s=0.25$ (reference) & 8 & 21 & 0 & 0 & 0 & 0 \\
\bottomrule
\end{tabular}
\end{table}

\begin{table}[!htbp]
\centering\small
\caption{Cheapest switching time without an added failure for each instance, chosen in hindsight, and the resulting mean cost and saving against the reference's 21 forward passes. Unlike Table~\ref{tab:fixed-fail}, each input gets its own switching time, so there are no added failures by construction; the choice needs the outcomes in advance and only bounds what an adaptive rule can save.}
\label{tab:fixed-oracle}
\begin{tabular}{@{}l r r r r r r@{}}
\toprule
\textbf{Prompts} & \textbf{Step 4} & \textbf{Step 5} & \textbf{Step 6} & \textbf{Step 8} & \textbf{Mean passes} & \textbf{Saving} \\
\midrule
Clean (500) & 488 & 9 & 3 & 0 & 13.06 & 37.8\% \\
Noisy (500) & 178 & 91 & 52 & 179 & 16.64 & 20.8\% \\
All (1{,}000) & 666 & 100 & 55 & 179 & 14.85 & 29.3\% \\
\bottomrule
\end{tabular}
\end{table}

 This per-input choice cannot be deployed, but it bounds what an adaptive rule can save, and the variation between inputs is the observation that motivates this study. Because the calibration and test splits enter this analysis and had already been inspected in earlier analyses, these figures are development evidence, and an independent evaluation needs data not used here. An earlier CPU version on 12 speakers (48 instances), with lighter evaluation models and a stricter failure rule, showed the same pattern.

\subsection{Student Restart Level}

Restarting the student one grid point later ($j=2$) saves one forward pass but produced more failures than the released code's default $j=1$ at every switching time tested: 26 versus 11 of 48 at $s=0.07$, 13 versus 8 at $s=0.10$, 9 versus 8 at $s=0.15$, and 14 versus 0 at $s=0.25$. The restart level is therefore fixed at $j=1$ (Section~\ref{sec:baseline}), and the switching time is the only handover setting the study varies.

\subsection{Stabilization Policy with a Cross-Validated Threshold}
\label{app:prelim-cv}

Unlike the other experiments in this appendix, this one was re-run on the Tesla T4 GPU with the evaluation models of Section~\ref{sec:measures}, on the 100 development speakers under the clean and the noisy15 prompt conditions (200 instances; \nolinkurl{code/analysis/cv_stabilization_t4.py}). The policy of Definition~\ref{def:policy}, with an SNR gate of 20~dB, was evaluated by four-fold speaker-grouped cross-validation. In each fold the threshold $\lambda$ was chosen on 75 speakers (150 instances) as the largest value in $\{0.20,0.25,\dots,0.60\}$ with zero added failures and frozen before scoring the 25 held-out speakers. The folds selected $\lambda=0.35$, 0.35, 0.35 and 0.50. Out of fold, mean cost fell from 21 to 18.88 forward passes ($-10.1\%$) with one added failure in 200 instances, against 50 of 200 for a fixed $s=0.15$, and the SNR gate sent all 100 noisy prompts to the reference, so the saving came from the clean prompts. Tables~\ref{tab:cv-policies} and~\ref{tab:cv-quality} compare it on the same 200 instances with the methods of Chapter~\ref{sec:method} (Table~\ref{tab:methods}), measured and tested as planned there. Against every fixed early switching time the policy had far fewer added failures (exact McNemar $p<10^{-13}$ for each), and the upper end of its 95\% interval for the failure rate, 2.8\%, stayed below the 5\% tolerance. Against the reference it used significantly fewer forward passes (Wilcoxon signed-rank test on speaker-level values, Holm-adjusted $p<10^{-15}$), while its differences in WER, SIM and UTMOS were not significant (Holm-adjusted $p=1.00$ for each); a claim that its quality equals the reference's needs the two one-sided tests with margins fixed before the held-out test. On clean prompts it saved 20.2\% of the passes and chose Steps~4, 5 and~6 for 6, 28 and 52 inputs, so the choice adapted per input. The SNR-gated controls behaved similarly on these data: SNR-only, the SNR gate followed by Step~5 or Step~6, and a random assignment of the policy's own steps among the gated inputs (2{,}000 draws, mean 0.75 failures) all stayed within 0 to 3 failures, and none differed significantly from the policy (exact McNemar $p\ge0.50$). Stabilization alone, without the gate, found no threshold with zero training failures in any fold and kept every input at the reference, because noisy prompts stabilize early (Appendix~\ref{app:candidates}). With at most three failures per method, these data cannot separate the SNR-gated variants; the comparison on the test split (Section~\ref{sec:drift-plan}) is designed for that. The earlier CPU version of this experiment, on 12 speakers with lighter evaluation models, had selected the edge of the search range, $\lambda=0.60$, in every fold; on the T4 data the selected thresholds lie inside the range, and widening it to 1.20 leaves them unchanged. Two limitations remain. Selecting a threshold with zero training failures is not a risk certificate: with 75 training speakers and no failures, the smallest risk that can be certified at $\delta=0.10$ is about 3\% (Appendix~\ref{app:calc-ltt}). And the development split also served other design decisions, so these out-of-fold figures may be optimistic. Objective~3 replaces this selection with Learn-Then-Test on speaker-disjoint data with one utterance per speaker.

\begin{table}[!htbp]
\centering\small
\caption{Compared methods of Chapter~\ref{sec:method} on the same 200 T4 development instances (100 speakers, clean and noisy15 prompts; \nolinkurl{code/analysis/cv_stabilization_t4.py}): added failures and cost. Stabilization only and the full policy are scored out of fold; cost-matched random is the mean over 2{,}000 random assignments of the full policy's steps among the gated inputs; the oracle needs the outcomes in advance and is shown as a bound. Upper 95\% is the upper end of the exact two-sided 95\% interval for the failure rate over all 200 instances \parencite{clopper1934}; McNemar is the exact test against the full policy.}
\label{tab:cv-policies}
\setlength{\tabcolsep}{4pt}
\begin{tabular}{@{}l r r r r r r@{}}
\toprule
& \multicolumn{2}{c}{\textbf{Failures}} & & & & \\
\cmidrule(lr){2-3}
\textbf{Method} & \textbf{Clean} & \textbf{Noisy} & \textbf{Upper 95\%} & \textbf{Passes} & \textbf{Saving} & \textbf{McNemar} \\
\midrule
Fixed $s=0.07$ (Step 4) & 3 & 63 & 40.0\% & 13.00 & 38.1\% & $<10^{-13}$ \\
Fixed $s=0.10$ (Step 5) & 2 & 49 & 32.1\% & 15.00 & 28.6\% & $<10^{-13}$ \\
Fixed $s=0.15$ (Step 6) & 0 & 50 & 31.6\% & 17.00 & 19.0\% & $<10^{-13}$ \\
Fixed $s=0.25$ (reference) & 0 & 0 & 1.8\% & 21.00 & 0.0\% & -- \\
SNR-only (Step 4) & 3 & 0 & 4.3\% & 17.00 & 19.0\% & 0.50 \\
SNR + fixed Step 5 & 2 & 0 & 3.6\% & 18.00 & 14.3\% & 1.00 \\
SNR + fixed Step 6 & 0 & 0 & 1.8\% & 19.00 & 9.5\% & 1.00 \\
Stabilization only & 0 & 0 & 1.8\% & 21.00 & 0.0\% & 1.00 \\
\textbf{Full policy} & 1 & 0 & 2.8\% & 18.88 & 10.1\% & -- \\
Cost-matched random & 0.75 & 0 & -- & 18.88 & 10.1\% & -- \\
\midrule
Oracle (hindsight) & 0 & 0 & 1.8\% & 14.70 & 30.0\% & -- \\
\bottomrule
\end{tabular}
\end{table}

\begin{table}[!htbp]
\centering\small
\caption{Mean quality of the compared methods on the instances of Table~\ref{tab:cv-policies}: WER with Whisper-large-v3, SIM with the WavLM-Base+ speaker verifier, and UTMOS. Over all instances, the mean WER of the fixed early switching times is dominated by a few noisy-prompt transcripts with many errors.}
\label{tab:cv-quality}
\setlength{\tabcolsep}{4pt}
\begin{tabular}{@{}l r r r r r r@{}}
\toprule
& \multicolumn{3}{c}{\textbf{Clean (100)}} & \multicolumn{3}{c}{\textbf{All (200)}} \\
\cmidrule(lr){2-4}\cmidrule(lr){5-7}
\textbf{Method} & \textbf{WER} & \textbf{SIM} & \textbf{UTMOS} & \textbf{WER} & \textbf{SIM} & \textbf{UTMOS} \\
\midrule
Fixed $s=0.07$ (Step 4) & 2.04\% & 0.934 & 3.86 & 34.93\% & 0.814 & 2.55 \\
Fixed $s=0.10$ (Step 5) & 2.08\% & 0.934 & 3.87 & 39.00\% & 0.816 & 2.56 \\
Fixed $s=0.15$ (Step 6) & 2.17\% & 0.934 & 3.86 & 8.78\% & 0.817 & 2.56 \\
Fixed $s=0.25$ (reference) & 2.27\% & 0.934 & 3.87 & 10.03\% & 0.821 & 2.56 \\
SNR-only (Step 4) & 2.04\% & 0.934 & 3.86 & 9.91\% & 0.821 & 2.55 \\
SNR + fixed Step 5 & 2.08\% & 0.934 & 3.87 & 9.93\% & 0.821 & 2.56 \\
SNR + fixed Step 6 & 2.17\% & 0.934 & 3.86 & 9.98\% & 0.821 & 2.56 \\
Stabilization only & 2.27\% & 0.934 & 3.87 & 10.03\% & 0.821 & 2.56 \\
\textbf{Full policy} & 2.13\% & 0.934 & 3.88 & 9.95\% & 0.820 & 2.56 \\
\midrule
Oracle (hindsight) & 2.04\% & 0.934 & 3.86 & 6.74\% & 0.828 & 2.55 \\
\bottomrule
\end{tabular}
\end{table}

The full policy was also timed with real early stopping on the Tesla T4 (\nolinkurl{research/evidence/benchmarks/gpu_timing_summary_by_condition.csv}), on the 200 test-split speakers under both prompt conditions, with GPU synchronization around a timed region that covered the prompt SNR estimate, the teacher steps with the stopping decision, the student and the vocoder (Table~\ref{tab:cv-runtime}). On clean prompts the mean synthesis time fell from 3.29~s at the reference to 2.40~s, 27\% less, which matches the 27\% fewer forward passes (15.30 instead of 21), so on this pipeline the pass saving carried over to time. Noisy prompts took 3.37~s under all three policies, so evaluating the gate and the stabilization score added no measurable time. Over both conditions the policy saved 0.44~s per utterance (95\% cluster-bootstrap interval 0.42 to 0.47~s). This run used $\lambda=0.50$, a value chosen from the earlier CPU labels rather than calibrated on the T4, and its speakers had been inspected in earlier analyses, so it is development evidence for the cost side of the method and not a test of its risk.

\begin{table}[!htbp]
\centering\small
\caption{Measured synthesis time per utterance on the Tesla T4 with real early stopping (200 test-split speakers; mean time and real-time factor).}
\label{tab:cv-runtime}
\setlength{\tabcolsep}{5pt}
\begin{tabular}{@{}l r r r r r@{}}
\toprule
& \multicolumn{3}{c}{\textbf{Clean}} & \multicolumn{2}{c}{\textbf{Noisy}} \\
\cmidrule(lr){2-4}\cmidrule(lr){5-6}
\textbf{Policy} & \textbf{Passes} & \textbf{Time} & \textbf{RTF} & \textbf{Time} & \textbf{RTF} \\
\midrule
Reference (Step 8) & 21.00 & 3.29~s & 0.73 & 3.37~s & 0.71 \\
SNR-only (Step 4) & 13.08 & 2.06~s & 0.45 & 3.37~s & 0.71 \\
Full policy ($\lambda=0.50$) & 15.30 & 2.40~s & 0.53 & 3.37~s & 0.71 \\
\bottomrule
\end{tabular}
\end{table}

\emph{Verdict and trade-off.} Taken together, these results support per-input switching and show where its trade-off lies. Choosing the switching time per input removed almost all of the failures that a fixed early switch causes (1 against 50 to 66 of 200) while still saving 10.1\% of the forward passes overall, 20.2\% on clean prompts and 27\% of the measured synthesis time on clean prompts, with no significant change in WER, SIM or UTMOS. The trade-off is between saving and caution: SNR-only saved more (19.0\%) with three failures, the full policy saved less (10.1\%) with one, and both stayed within the 5\% tolerance at the upper end of their intervals. On these data the stabilization score made the policy more cautious, but its step choices could not yet be distinguished from the SNR-gated controls, and noisy prompts kept the reference cost because the gate sends them there. The calibration of Objective~3 will replace the cross-validated threshold with a certified one, and the held-out comparison of Objective~4 will decide which variant to recommend.

\subsection{Development-Split Signal Diagnostic}
\label{app:prelim-dev}

A first diagnostic of the two signals was run on the 100 development speakers with the evaluation models of this study (Section~\ref{sec:measures}), each under the clean and the 15~dB noisy condition. The estimated prompt SNR separated the two conditions completely: 23.3 to 50.0~dB for clean prompts and 16.8 to 18.2~dB for noisy prompts, so a 20~dB gate sent every noisy prompt to the reference. Fixed early handovers failed on 49\% to 63\% of noisy prompts but on only 2\% to 4\% of clean ones. On clean prompts, the stabilization score $d_4$ did not separate Step-4 failures from passes (area under the receiver operating characteristic (ROC) curve about 0.47, based on only two to three failures), and at $\lambda=0.50$ the policy waited past Step~4 on 78 of the 100 clean prompts, of which 75 to 77 would already have passed at Step~4. These figures come from the development split alone, which serves design decisions, and the small number of clean failures makes the estimates imprecise. They motivate the comparisons of Appendix~\ref{app:drift-plan} but are not evidence for or against the stabilization score.

\subsection{Plan for Further Investigation of the Stabilization Signal}
\label{app:drift-plan}

Because the value of the stabilization score is uncertain, the implementation after the proposal will test it in four steps, each fixed before the data it uses are examined. The second step is the comparison summarized in Section~\ref{sec:drift-plan}.
\begin{enumerate}[leftmargin=1.4em]
\item \textbf{Discrimination.} On the development split, the study will measure how well the score separates handovers that fail from handovers that pass, using the area under the ROC curve, rank correlation with the per-utterance quality change, and the variance it explains beyond the prompt SNR.
\item \textbf{Controls.} On the test split, the policy that uses both signals will be compared with the SNR gate followed by a fixed later step and with a random choice of steps that matches its mean cost, using the exact McNemar test on paired failures. If the stabilization score does not beat these controls, its step choices carry no information beyond waiting longer.
\item \textbf{Conditions the SNR cannot see.} On speakers not used in any split, the comparison will be repeated on prompts degraded in ways that the energy-based SNR estimate may miss, such as reverberation and telephone-band recordings, with the hypothesis and test written down before synthesis.
\item \textbf{Richer zero-cost signals.} Variants of the score that need no extra forward pass, such as per-frame or per-band changes and the gap between the conditional and unconditional predictions, will be explored on the development split only; any variant that is kept will be calibrated and tested on data it has not seen.
\end{enumerate}
If the stabilization score adds no measurable value in these tests, the SNR-gated variant of the policy will be recommended, and the stabilization score will be reported as a tested negative result together with the reasons found for it. The contribution of the study, a per-utterance handover framework with a certified bound on added failures, does not depend on which signal performs best.

\begin{table}[!htbp]
\centering\small
\caption{Preliminary experiments and their consequences for the design.}
\label{tab:prelim}
\begin{tabularx}{\textwidth}{@{}>{\raggedright\arraybackslash}p{3.0cm} >{\raggedright\arraybackslash}X >{\raggedright\arraybackslash}X@{}}
\toprule
\textbf{Experiment} & \textbf{Finding} & \textbf{Consequence for this study} \\
\midrule
Endpoint cushioning & Partial acceptance of the student's first correction gave no detectable benefit & Not pursued; the handover itself is left unchanged \\
Fixed switching times & On 1{,}000 T4 instances, the failure rate rises as the switch moves earlier (25\% to 33\%); the cheapest acceptable switch varies by input; a per-input choice in hindsight needs 14.85 of 21 passes & Motivates per-utterance switching (Objective~1) \\
Student restart level & Restarting at $j=2$ adds failures at every switching time & Restart fixed at $j=1$; switching time is the only variable \\
Cross-validated stabilization policy & On T4 development data, 10\% fewer passes with one added failure in 200 out of fold, and 27\% less synthesis time on clean prompts with real early stopping, but no risk certificate & Motivates the decision signals (Objective~2) and calibration with Learn-Then-Test (Objective~3) \\
\bottomrule
\end{tabularx}
\end{table}

\FloatBarrier


% =====================================================================
\section[APPENDIX E.\ CANDIDATE SOLUTIONS]{Candidate Solutions}
% =====================================================================
\label{app:candidates}

The preliminary experiments of Appendix~\ref{app:prelim} were followed by four development studies on Google Colab GPUs, which tested gates, handover operators and noise-aware variants of the sampler. Two candidates remain. They extend, and do not replace, the method of Chapter~\ref{sec:method}: the \emph{Trajectory Anomaly Score} (TAS) decides which inputs must not be handed over early, and \emph{SNR-Wiener Conditioning} (SWC) removes the prompt's noise so that more inputs can be handed over early. Both use the prompt SNR estimate of Definition~\ref{def:snr} and add no network evaluation. They are reported here as development evidence, and each will be tested on speakers not used here before it is adopted (Section~\ref{app:cand-limits}).

The studies used the first 60 development speakers, each under six prompt conditions (the clean prompt and white Gaussian noise added at 15, 18, 20, 22 and 25~dB), giving 360 instances scored with the evaluation models of Section~\ref{sec:measures} and the added-failure rule of Definition~\ref{def:failure} against the reference switching time. Every threshold was certified by the fixed-sequence procedure of Algorithm~\ref{alg:ltt} ($\alpha=0.05$, $\delta=0.10$) on the training speakers of a five-fold speaker-grouped cross-validation, and all figures below are out of fold. The comparison baseline is the SNR-only policy: a handover at teacher step~4 (13 forward passes) if $\hat S\ge20$~dB, and the reference (21 forward passes) otherwise. On these instances it failed on 86 of 360 (23.9\%), because prompts with 18 to 22~dB of added noise are estimated above 20~dB and pass the gate.

\subsection{Trajectory Anomaly Score}
\label{app:cand-tas}

\emph{Motivation.} By the endpoint formula of \Cref{eq:endpoint} and the Euler step of \Cref{eq:euler}, consecutive endpoint estimates differ by
\[
\hat x_1(t_i)-\hat x_1(t_{i-1}) = (1-t_i)\bigl(v(x_{t_i},t_i)-v(x_{t_{i-1}},t_{i-1})\bigr),
\]
so the stabilization score $d_i$ of Definition~\ref{def:stab} measures how much the generation direction is still changing. A noisy prompt weakens the conditioning: the median $d_4$ fell from 0.61 on clean prompts to 0.17 at 15~dB, together with the relative gap between the conditional and unconditional predictions (median 0.236 to 0.116). The trajectory therefore settles early because the condition carries less information, not because the output is ready. Under this \emph{premature stabilization}, the stopping rule of Definition~\ref{def:policy} hands noisy prompts over early, and on these instances no threshold of that rule could be certified.

\emph{Definition.} For an input $x$, let $\mathbf z(x)=(\log d_1,\dots,\log d_6)$ collect the stabilization scores that the teacher pass already computes. From the clean prompts of the calibration speakers, without failure labels, estimate their mean $\boldsymbol\mu$ and covariance $\boldsymbol\Sigma$. The anomaly score counts only trajectories that are stiller than normal:
\[
A(x)=\begin{cases}(\mathbf z-\boldsymbol\mu)^{\top}\boldsymbol\Sigma^{-1}(\mathbf z-\boldsymbol\mu) & \text{if the mean component of } \mathbf z-\boldsymbol\mu \text{ is negative},\\[2pt] 0 & \text{otherwise.}\end{cases}
\]
The handover rule keeps the SNR gate and adds the score:
\[
\pi_{\mathrm{TAS}}(x)=\begin{cases}4 & \text{if } \hat S\ge\theta \text{ and } A(x)\le\hat\kappa,\\[2pt] k_{\text{ref}} & \text{otherwise,}\end{cases}
\]
where $\theta$ is the SNR gate (20~dB in these studies, as in SNR-only) and $\hat\kappa$ is the most permissive threshold certified by Algorithm~\ref{alg:ltt}, testing thresholds in increasing order.

\emph{Relation to prior work.} Mahalanobis distances are standard anomaly scores, and DiffPath uses the rate of change and curvature of diffusion paths to detect out-of-distribution images \parencite{heng2024diffpath}. TAS differs in what it detects (corrupted conditioning), in treating a trajectory that settles too early as a reason not to stop, the opposite of early-exit rules that stop once the output stops changing \parencite{schuster2022calm}, and in gating a calibrated teacher-to-student handover. To our knowledge this use has not been reported.

\subsection{SNR-Wiener Conditioning}
\label{app:cand-swc}

\emph{Motivation.} On the noisy prompts most failures were speaker-similarity failures: the generated voice inherits the prompt's noise. Changing only what the student receives at the handover, or steering the teacher's guidance, recovered at most 0.024 of speaker similarity, whereas the noise itself cost 0.06 to 0.24, measured on the same speakers and seeds with the clean and the noisy prompt.

\emph{Definition.} Before the teacher and the student read the prompt, the prompt frames of their conditioning are cleaned in the model's own mel domain. For mel band $b$, let $N_b$ be the mean linear magnitude of the quietest 20\% of the prompt frames, and let $m_{f,b}$ be the magnitude of prompt frame $f$. Each prompt frame is multiplied by
\begin{gather*}
G_{f,b}=\min\Bigl\{1,\ \max\Bigl\{0.1,\ 1-\beta(\hat S)\,\frac{N_b}{m_{f,b}}\Bigr\}\Bigr\},\\
\beta(\hat S)=\min\Bigl\{1,\ \max\Bigl\{0,\ \frac{30-\hat S}{10}\Bigr\}\Bigr\},
\end{gather*}
so that prompts estimated at 30~dB or above are left unchanged and prompts at 20~dB or below are cleaned at full strength. The teacher then runs to step~4 and the student completes as usual, for 13 forward passes.

\emph{Routing.} The prompt SNR also decides which inputs are cleaned. With two thresholds $\theta_{\mathrm{e}}\le\theta_{\mathrm{hi}}$, certified in turn by Algorithm~\ref{alg:ltt},
\[
\pi_{\mathrm{SWC}}(x)=\begin{cases}4 \text{ with the original prompt} & \text{if } \hat S\ge\theta_{\mathrm{hi}},\\[2pt] 4 \text{ with the cleaned prompt} & \text{if } \hat S<\theta_{\mathrm{e}},\\[2pt] k_{\text{ref}} & \text{otherwise.}\end{cases}
\]
Setting $\theta_{\mathrm{e}}=\theta_{\mathrm{hi}}$ cleans every input below the gate; a lower $\theta_{\mathrm{e}}$ cleans heavy noise only and sends moderately noisy prompts to the reference.

\emph{Relation to prior work.} The gain is classical spectral subtraction with a minimum-energy noise estimate \parencite{boll1979spectral}, so SWC is not a new formula. What is new, to our knowledge, is applying it to the conditioning of a teacher-guided sampler, with the strength set by the prompt SNR, as a substitute for teacher computation that makes an early handover safe.

\subsection{Results}
\label{app:cand-results}

Table~\ref{tab:candidates} compares the candidates with the SNR-only policy on the same 360 instances. Both candidates removed most of SNR-only's failures, which fell mainly on the moderately noisy prompts, and SWC on heavy noise did so while saving 19.1\% of the forward passes. The thresholds were certified per fold on the training speakers; the pooled out-of-fold rates of TAS and of SWC on heavy noise (3.3\% and 3.1\%) lie below the 5\% tolerance, while the upper ends of their 95\% intervals over 360 instances (5.8\% and 5.4\%) lie slightly above it, which the larger calibration and test splits of Chapter~\ref{sec:method} would narrow. On the 300 noisy instances, SWC's step-4 output was better than the reference output on average: speaker similarity rose by 0.126 (higher in 90\% of the instances; Wilcoxon signed-rank $p<10^{-40}$) and WER fell by 4.8 percentage points. With every input handed over at step~4, SWC failed on 21 of the 300 noisy instances, against 126 for the plain step-4 handover. With $\theta_{\mathrm{e}}=\theta_{\mathrm{hi}}$ (26~dB in every fold), every input below the gate was cleaned, and the policy had fewer failures and fewer forward passes than SNR-only; with heavy noise only ($\theta_{\mathrm{e}}$ between 18 and 21~dB), it kept the failure rate within the 5\% tolerance.

TAS also kept the failure rate within the tolerance that SNR-only exceeds. It did not improve on a fixed SNR gate whose threshold was itself calibrated on these noise levels (7 failures at 18.71 forward passes; the difference from TAS is not significant, exact McNemar $p=0.23$), but it needs only clean prompts to model normal trajectories. The stopping rule of Definition~\ref{def:policy} certified no threshold on these instances and fell back to the reference for every input.

Other variants did not improve on the plain handover: guidance steered away from a noise-only prompt, a guidance scale raised with the SNR, an antithetic pair of student runs, and noise cancellation applied to the generated output instead of the prompt. A pretrained speech-enhancement network applied to the prompt helped about as much as SWC (speaker similarity $+0.117$) but adds a network to the pipeline.

\begin{table}[!htbp]
\centering\small
\caption{Candidate policies against SNR-only on 360 development instances (60 speakers; clean, and five noisy conditions of 60 instances each; out of fold, $\alpha=0.05$, $\delta=0.10$). Upper 95\% is the upper end of the exact two-sided 95\% interval for the failure rate over all 360 instances \parencite{clopper1934}; McNemar is the exact test against SNR-only. Of SNR-only's failures, TAS avoids 74 and adds none, SWC with $\theta_{\mathrm{e}}=\theta_{\mathrm{hi}}$ avoids 75 and adds 9, SWC on heavy noise avoids 79 and adds 4, and the calibrated SNR gate avoids 79 and adds none.}
\label{tab:candidates}
\setlength{\tabcolsep}{4pt}
\begin{tabularx}{\textwidth}{@{}>{\raggedright\arraybackslash}X r r r r r r@{}}
\toprule
& \multicolumn{2}{c}{\textbf{Failures}} & & & & \\
\cmidrule(lr){2-3}
\textbf{Policy} & \textbf{Clean} & \textbf{Noisy} & \textbf{Upper 95\%} & \textbf{Passes} & \textbf{Saving} & \textbf{McNemar} \\
\midrule
SNR-only ($\hat S\ge20$~dB) & 2 & 84 & 28.6\% & 14.40 & 31.4\% & -- \\
Calibrated SNR gate & 2 & 5 & 4.0\% & 18.71 & 10.9\% & $3{\times}10^{-24}$ \\
\textbf{SNR + TAS} & 2 & 10 & 5.8\% & 18.20 & 13.3\% & $1{\times}10^{-22}$ \\
\textbf{SNR + SWC}, $\theta_{\mathrm{e}}=\theta_{\mathrm{hi}}$ & 2 & 18 & 8.4\% & 13.00 & 38.1\% & $4{\times}10^{-14}$ \\
\textbf{SNR + SWC}, heavy noise & 2 & 9 & 5.4\% & 16.98 & 19.1\% & $4{\times}10^{-19}$ \\
Reference for every input & 0 & 0 & 1.0\% & 21.00 & 0.0\% & $3{\times}10^{-26}$ \\
\bottomrule
\end{tabularx}
\end{table}

\subsection{Limitations and Next Steps}
\label{app:cand-limits}

These results are development evidence. TAS was designed after premature stabilization was observed in these same instances, the added noise was white Gaussian noise only, and 60 speakers are fewer than the calibration and test splits of Chapter~\ref{sec:method}. The SWC outputs were synthesized on an A100 GPU and compared with reference outputs from the Tesla T4: the teacher reproduced exactly, but the student's random draws differ between the two GPUs, so each SWC output was judged against a reference made with different student noise. This adds disagreements that inflate the failure counts of SWC, while those of TAS and SNR-only, which use T4 outputs only, are unaffected. Before either candidate is adopted, the study will repeat SWC with its references on the same GPU, including an SWC output at the reference switching time to separate the gain from cleaning from the safety of the early handover; test both candidates, with their settings fixed in advance, on speakers not used here; and test prompts with babble, reverberation and telephone-band degradation, where the SNR estimate is less reliable and SWC's assumption of stationary noise is weakest.

\FloatBarrier
